% Ad M. Brutum 1.6 (ad-brutum-01-06) --- English translation
% Translated by Claude (Anthropic) from the Latin (Perseus).
% Style guide: STYLE.md.

\ciceroLetterOpener{Brutus Ciceroni salutem}

\ciceroSection{1} Do not wait for me to thank you. From the kind of bond between us --- which has reached the height of goodwill --- such an exchange ought long since to have been done away with. Your son is away from me; we shall meet in Macedonia. He has been ordered to lead the cavalry from Ambracia through Thessaly. I have written to him to come to meet me at Heraclea. When I have seen him, since you grant us leave, we shall decide jointly about his return for standing for office or for being put up for one.

\ciceroSection{2} I commend to you in the warmest terms Glyco, the physician of Pansa, who has the sister of our friend Achilleos in marriage. We hear that he has come under suspicion with Torquatus over the death of Pansa, and is being held in custody as a parricide. Nothing could be less to be believed; for who has taken a greater blow from Pansa's death than he? Besides, he is a modest and honest man, whom not even self-interest, by the look of it, would have driven to such a crime. I beg you --- and I beg you in earnest, for our friend Achilleos is no less distressed than he has every right to be --- rescue him from custody and save his life. I count this as belonging to my duty in private matters as much as anything else does.

\ciceroSection{3} While I was writing this letter to you, a dispatch was delivered to me by Satrius, the lieutenant of C. Trebonius, brought by Tillius and Deiotarus, that Dolabella has been cut down and put to flight. I am sending you a Greek letter from a certain Cicereius, sent to Satrius.

\ciceroSection{4} Our friend Flavius has, in connection with a dispute about an inheritance which he has with the people of Dyrrachium, taken you as his judge. I ask you, Cicero, and Flavius asks too: please settle the matter. That the man who made Flavius his heir owed money to the community is not in doubt; nor do the people of Dyrrachium deny it, but they say the debt was forgiven them by Caesar. Do not allow my friend to suffer wrong at the hands of your friends. 19 May, from camp at the foot of Candavia.
